\section{Methodology}

This study proposes a three-stage approach to assess the impact of fossil fuel price fluctuations on WACC for energy transition projects. The end-to-end tensor workflow is summarized in Fig.~\ref{fig:tensor-framework}, and the stage-specific model relations are detailed in Fig.~\ref{fig:methodology-model-relations}. The first stage predicts fossil fuel prices using a hybrid LSTM-GPR model, and the second stage calculates the corresponding inflation rates.

After that, the third stage evaluates the impact of fossil fuel prices on WACC using an AutoGluon model based on performance metrics of $R^2$, MAPE, and RMSE.

\begin{figure*}[t]
    \centering
    \PaperMethodology{0.98\textwidth}{fossilfuel_wacc_methodology_model_relations.drawio.pdf}
    \caption{Methodology and model-relation overview for the three-stage fossil-fuel price, inflation, and WACC assessment pipeline.}
    \label{fig:methodology-model-relations}
\end{figure*}

\subsection{Stage 1: Prediction of Fossil Fuel Prices}

Since previous studies have proved the proficiency of Long Short-Term Memory (LSTM) in handling sequence data with long-term dependencies \cite{ALSELWI2024102068,LINDEMANN2021650}, and evidence on Gaussian Process Regression (GPR) for fossil-energy forecasting \cite{LAHMIRI2024105008}, together with its probabilistic formulation \cite{rasmussen2003gaussian}, this study uses a hybrid model combining LSTM and GPR to predict fossil fuel prices (coal, oil, and natural gas) from 2026 to 2035 (Fig.~\ref{fig:lstm}).

\subsubsection{Step 1.1: LSTM Model}
The primary forecasting model is an LSTM network designed to learn intricate non-linear patterns from historical time-series data. To represent the cyclical nature of monthly patterns, we transform the month of the year, $m$, into two continuous sinusoidal features:
\begin{equation}
    m_{\sin}=\sin\left(\frac{2\pi m}{12}\right),\quad
    m_{\cos}=\cos\left(\frac{2\pi m}{12}\right).
\end{equation}
The model architecture integrates a convolutional layer for local pattern extraction, followed by bidirectional LSTM and GRU layers to capture long-range temporal dependencies. An attention mechanism is subsequently applied, enabling the model to dynamically weight the importance of different time steps in the input sequence during prediction generation. The model predicts the price for the next time step, and this prediction is then used to update the feature set for the subsequent step, which is fed back into the model to generate the next forecast.

\subsubsection{Step 1.2: Forecast Refinement with Gaussian Process Regression}
The point estimates generated by the LSTM model are subsequently refined using a GPR model. We employ a composite kernel to model the underlying price function:
\begin{equation}
    k(t,t')=C\cdot\exp\left(-\frac{\|t-t'\|^2}{2l^2}\right)+\sigma_n^2\delta(t,t').
\end{equation}
where $C$ is the constant kernel, $l$ is the length-scale parameter, and $\sigma_n^2\delta(t,t')$ represents the white-noise regularization term.

This kernel incorporates a Radial Basis Function component that encodes the assumption of smoothness in the underlying function space. Given the initial predictions from the LSTM model, represented as the vector $\mathbf{y}_{DL}$, the GPR computes a predictive mean, serving as the price $P_e$ for Stage 3:
\begin{equation}
    P_e=K(t,X)[K(X,X)+\sigma_n^2 I]^{-1}\mathbf{y}_{DL}.
\end{equation}
where $K(t,X)$ is the vector of covariances between the test point $t$ and the training inputs $X$, and $K(X,X)$ is the covariance matrix of the training inputs. This GPR stage smooths the neural forecast path; its conditional variance does not capture the full uncertainty of the preceding neural forecast. The GPR component uses 10 optimizer restarts, a regularization value $\alpha=0.01$, and a fixed random seed of 42. The main LSTM and GPR settings are summarized in Appendix Table~\ref{tab:hyperparameter-summary}.

\begin{figure}[t]
    \centering
    \PaperNetwork{0.92\columnwidth}{LSTM.png}
    \caption{Architecture for the price-prediction model. It includes a network with CNN, LSTM/GRU, attention (A), and normalization (N) layers, which is then smoothed by GPR.}
    \label{fig:lstm}
\end{figure}

\subsection{Stage 2: Calculation of Inflation Rate}

\subsubsection{Inclusion of Inflation Rate as a Covariate}
In the model, the inflation rate is included as a crucial covariate. As a macroeconomically sensitive indicator, WACC is directly influenced by the prevailing interest rate environment, and inflation is relevant to monetary-policy and financing conditions \cite{dobrowolski2022weighted}. Additionally, fossil fuel prices themselves are closely correlated with macroeconomic inflation \cite{millischer2024renewables}.

\subsubsection{Forecasting Framework}
This study adopts an inflation-forecasting architecture informed by the recurrent-network approach of Paranhos \cite{PARANHOS2025}, which models the inflation rate as a non-linear function fitting problem:
\begin{equation}
y_{t+h}=G(x_t;\Theta_h)+\epsilon_{t+h}.
\end{equation}
where $y_{t+h}$ is the inflation rate at a future horizon $h$, and $x_t$ is the vector of predictor variables containing relevant macroeconomic indicators and their lags.

The function $G(\cdot)$ combines an LSTM structure with a fully connected feed-forward (FF) network:
\begin{equation}
G(x_t;\Theta_h)=g_{FF}(\mathrm{LSTM}(x_t;\Theta_{LSTM});\Theta_{FF}).
\end{equation}
In this architecture, the LSTM first processes the sequential input data to extract a set of dynamic factors. Subsequently, the feed-forward network takes these dynamic factors as input to generate the final inflation forecast.

To address the high sensitivity of neural network training to random parameter initializations, we employ an ensemble learning strategy for model optimization. Specifically, the final forecast is calculated by averaging the outputs of $N=5$ independently trained models, each with a different initialization:
\begin{equation}
\hat{y}_{ens,t+h}=\frac{1}{N}\sum_{n=1}^{N}\hat{y}_{n,t+h}.
\end{equation}
where $\hat{y}_{n,t+h}$ is the prediction from the $n$-th model.

The parameters $\Theta_h$ for all models are estimated by minimizing MSE:
\begin{equation}
\hat{\Theta}_{h}=\arg\min_{\Theta_h}\left\{\frac{1}{T-h}\sum_{t=1}^{T-h}(y_{t+h}-G(x_t;\Theta_h))^2\right\}.
\end{equation}

\subsection{Stage 3: Evaluation of WACC Based on Fossil Fuel Price}

To quantify the impact of fossil fuel price fluctuations on the WACC, we employed an Automated Machine Learning (AutoML) model, which can automate model selection, hyperparameter tuning, and ensembling.

\subsubsection{Modeling Design}
This assessment model involves three fossil fuel prices (coal, natural gas, and oil) and nine types of energy transition projects. For each region, we constructed $3\times9=27$ separate regression specifications, giving 81 region--fuel--project combinations across the three regions. Coal and gas generation are retained as conventional-technology comparisons. The target WACC series are scenario projections derived from Calcaterra et al. \cite{calcaterra2024reducing}, aligned to 120 monthly rows per region for 2026--2035. In each model, the independent variables consist of an energy price and the contemporaneous inflation rate, while the dependent variable is the WACC for the project. This approach aims to find a function $f$ for each project type $p$ and energy source $e$, such that:
\begin{equation}
    \mathrm{WACC}_{p}=f(P_e,\hat{y}_{ens,t})+\epsilon.
\end{equation}
where $\mathrm{WACC}_{p}$ is the capital cost for a specific project, $P_e$ is the price, $\hat{y}_{ens,t}$ represents the inflation rate, and $\epsilon$ is the irreducible error term.

\subsubsection{Automated Modeling with AutoGluon}
We selected AutoGluon-Tabular for training and optimization \cite{autogluon2025docs}. The implementation uses two ensembling techniques:
\begin{itemize}
    \item \textbf{$j$-fold Cross-Validation Bagging}: It trains multiple base models $h_j$ on different subsets of the training data created via cross-validation. The final result, $\hat{y}_{bag}$, is the average of the output of all base models:
    \begin{equation}
    \hat{y}_{bag}(x)=\frac{1}{J}\sum_{j=1}^{J}h_j(x).
    \end{equation}
    \item \textbf{Auto Stacking}: Let the set of base models be $\{h_m\}_{m=1}^{M}$. A higher-level meta-model, $H$, is trained on the out-of-fold predictions of these base models. The final result, $\hat{y}_{stack}$, is generated by this meta-model, which learns the optimal way to combine the base model outputs:
    \begin{equation}
    \hat{y}_{stack}(x)=H(h_1(x),h_2(x),\ldots,h_M(x)).
    \end{equation}
\end{itemize}
Through this automated process, we generated an optimized model for each region-specific set of 27 regressions, best adapted to its data characteristics.

\subsubsection{Model Performance Evaluation}
The archived AutoGluon implementation evaluates fitted models on the same input rows used for training. Its $R^2$, MAPE, and RMSE therefore summarize in-sample fit. Here $R^2$ measures variance explained by the joint fuel-price and inflation specification; MAPE is reported as a fraction, and RMSE is in decimal WACC units. Fit statistics do not measure the size of a price response.

Additional diagnostics use cleaned panels with region-specific inflation. OLS and Ridge specifications are evaluated using a fixed chronological holdout (2026--2032 training; 2033--2035 testing) and rolling-origin folds. These checks examine the stability of the modeled relationships rather than serving as an out-of-sample evaluation of AutoGluon itself. Log-linear specifications control for inflation and a time trend. Slopes are transformed back from standardized inputs, giving the elasticity $\beta_e=\partial\log(\mathrm{WACC})/\partial\log(P_e)$. A 10\% price increase is converted into a WACC change as $10^4\overline{\mathrm{WACC}}[\exp(\beta_e\log(1.1))-1]$ basis points. Granger-style tests examine historical lagged macroeconomic associations.

\begin{figure*}[t]
    \centering
    \PaperOilComparison{0.92\textwidth}{Accuracy.png}
    \caption{Accuracy of the prediction model: forecast results of oil prices from 2021 to 2024 using different models (LSTM-GPR, SARIMA, XGBoost, LightGBM).}
    \label{fig:accuracy}
\end{figure*}

\subsection{Accuracy of the Prediction Models}
The archived oil-price comparison contains 47 monthly observations from 2021--2024 for LSTM-GPR, SARIMA, XGBoost, and LightGBM (Fig.~\ref{fig:accuracy}). Recomputing metrics from \texttt{Accuracy\_Prediction-Model.xlsx} gives $R^2=0.963$, MAPE=0.0186, RMSE=2.416, and MAE=1.425 for LSTM-GPR (Table~\ref{tab:stage1-accuracy}). The LSTM-GPR series has the smallest errors among these archived outputs. The saved comparison outputs lack complete training-origin metadata. Fixed-origin baseline forecasts trained through 2020 are therefore reported separately in the repository. The inflation model follows the recurrent-network approach of Paranhos \cite{PARANHOS2025}.

\begin{table}[t]
\centering
\caption{Archived Stage 1 Oil-Price Comparison, 2021--2024}
\label{tab:stage1-accuracy}
\scriptsize
\begin{tabular}{lrrrr}
\toprule
\textbf{Model} & \textbf{$R^2$} & \textbf{MAPE} & \textbf{RMSE} & \textbf{MAE} \\
\midrule
LSTM-GPR & 0.963 & 0.0186 & 2.416 & 1.425 \\
LightGBM & 0.932 & 0.0299 & 3.267 & 2.464 \\
XGBoost & 0.729 & 0.0644 & 6.546 & 5.273 \\
SARIMA & -10.206 & 0.4839 & 42.069 & 39.449 \\
\bottomrule
\end{tabular}
\end{table}
